\centerline{\bf \Huge Неравенство о средних. Начало.}

\begin{flushright}
        \scriptsize{}
\end{flushright}

В этом листочке мы научимся пользоваться следующим очень полезным неравенством

\[\dfrac{a_1 + a_2 + \ldots + a_n}{n} \geq \sqrt[n]{a_1a_2\ldots a_n}\]

\problem{0.1} \textit{Неравенство Коши.} 
Докажите, что для любых положительных $x, y$ верно неравенство 
\[x^2 + y^2 \geq 2xy\]

\textit{Следствие. Неравенство о средних, $n = 2.$} Пусть $x = \sqrt{a}, y = \sqrt{b}.$ Тогда верно, что
\[\dfrac{a+b}{2} \geq \sqrt{ab}.\]

\problem{0.2} \textbf{Неравенство о средних, $n=4$.} Докажите, что для любых положительных $a, b, c, d$ верно неравенство

\[\dfrac{a+b+c+d}{4} \geq \sqrt[4]{abcd}\]

\problem{0.3} \textit{Очевидный индуктивный переход.} Докажите неравенство о средних верно для любых степеней двойки.

\problem{0.4} Докажите, что если неравенство о средних верно для $n$, то оно верно и для $n-1$. Тем самым завершите доказательство неравенства о средних.

\problem{1} 
Для положительных $a, b, c$ докажите, что
$(a^2b + b^2c + c^2a)(ab^2 + bc^2 + ca^2) \geq 9a^2b^2c^2$

\problem{2}
Докажите, что для положительных \( x \) выполнено \[ x^{40} + \frac{1}{x^{16}} + \frac{2}{x^4} + \frac{4}{x^2} + \frac{8}{x} \geq 16.\]

\problem{3}
Положительные числа \( x, y, z \) таковы, что \( xy^2z^3 = 108 \). Какое минимальное значение может принимать \( x + y + z \)?

\newpage

\hspace{3mm}

\problem{4}
Докажите, что для положительных чисел $a_1, a_2, \ldots, a_n$ выполняется:

\[\dfrac{a_1+a_2+\ldots+a_n}{n} \leq \sqrt{\dfrac{a_1^2 + a_2^2 + \ldots + a_n^2}{n}}\]

\problem{5}
Сумма положительных чисел $a, b, c, d$ равна 1. Докажите, что 

\[\sqrt{1+4a}+\sqrt{1+4b}+\sqrt{1+4c}+\sqrt{1+4d} \leq 4\sqrt{2}\]

\problem{6}
Для положительных чисел \( a_1, a_2, \ldots, a_n \) докажите неравенство
\[
\sqrt[n]{a_1 a_2 \ldots a_n} \geq \frac{n}{\frac{1}{a_1} + \frac{1}{a_2} + \ldots + \frac{1}{a_n}}
\]

\problem{7}
Для положительных \( a, b, c \) докажите, что
\[
\frac{9}{a + b + c} \leq 2 \left( \frac{1}{a + b} + \frac{1}{b + c} + \frac{1}{a + c} \right)
\]

\problem{8}
$\bigl[\textbf{ДВИ в МГУ 2023}\bigr]$
Положительные числа $a,b,c$ таковы, что
\[a^2 + b^2 + c^2 = 1\]
Найдите наибольшее значение выражения $ab + bc\sqrt{3}$